% Continuum-bridge theory package, v3 after TWO adversarial verification rounds.
% Round-2 fixes: contour orientation phrasing; (H1') affine-theta clause; delta_0 gains the
% fourth term 5*sigma*r0/(8*sqrt2) (round-2 fatal, counterexample-backed); omega_J entrywise
% -> operator-norm factor 2 repaired via omega_J(r0) <= sigma/16; margin bookkeeping now
% direct from the N-wide hypotheses (3/4-machinery removed); converse fully hypothesized
% (fold existence, boundary distance, uniform moduli) with explicit uniform constant;
% r0 also capped by dist(q*, boundary N); moduli are data (omega), Lambda is sup-norm bound;
% R2 same-branch qualifier + "cancel to leading order"; R3 explicit geometric domination.

% ============ MAIN TEXT STATEMENTS ============

\subsection{A continuum bridge with a posteriori control}
\label{sec:bridge-theory}

The refinement study below is organized by two exact statements and one
conditioned remark. Write \(F_h(t,\theta)=f_h(t,\theta)\) for the level-\(h\)
reaction density and \(F\) for a candidate continuum limit on a compact
rectangle \(N=[t_1,t_2]\times[\theta_1,\theta_2]\) with \(t_1>0\).

\paragraph{Jet convergence from sectorial resolvent convergence.}
Let \(A_h(\theta)=L_h-\operatorname{diag}k_{h,\theta}\) be the killed
generators, with
\(f_h(t,\theta)=\langle k_{h,\theta},\e^{A_h(\theta)t}p_h\rangle\), and let
\(A(\theta)\) be the continuum killed operator. Assume:
(H1) the family \(\{A_h(\theta)\}\) is uniformly sectorial---there are
\(M\ge1\), \(\omega\in\mathbb R\), \(\delta\in(0,\pi/2)\), independent of
\(h\) and \(\theta\), with \(\|(z-A_h(\theta))^{-1}\|\le M/|z-\omega|\) for
\(|\arg(z-\omega)|<\tfrac\pi2+\delta\);
(H1$'$) \(A(\theta)\) is sectorial on the same sector with the analogous
bound, uniformly in \(\theta\); the continuum family is affine along the
same path, \(A(\theta)=A(0)-\theta\,(k_{\mathrm{pat}}-\bar k)\) acting by
multiplication, with \(\theta\)-differentiable resolvent; the pairing
\(\langle k_\theta,\cdot\,\rangle\) with \(p\) is well defined uniformly in
\(\theta\); and \(f(t,\theta)=\langle k_\theta,\e^{A(\theta)t}p\rangle\)
admits the analytic-semigroup representation below, so that
\(\partial_\theta\partial_t^jf\) has the same two-term Dunford form as its
discrete counterpart. Fix \(\delta'\in(0,\delta)\), \(\rho>0\), and the
\emph{unbounded} sectorial contour
\(\Gamma=\{\omega+r\e^{\pm i(\pi/2+\delta')}:r\ge\rho\}\cup
\{\omega+\rho\e^{i\varphi}:|\varphi|\le\tfrac\pi2+\delta'\}\),
oriented from the lower ray to the upper ray (counterclockwise about the
spectrum). For \(t\ge t_1>0\) the Dunford representation
\[
\partial_t^jf_h(t,\theta)
=\frac1{2\pi i}\int_\Gamma z^j\e^{zt}\,
\langle k_{h,\theta},(z-A_h(\theta))^{-1}p_h\rangle\,\dd z,
\qquad j\le3,
\]
holds exactly for every \(h\) (and, by (H1$'$), for the continuum pair).
Define the scalar contour defects
\[
\eta_h=\sup_{\theta}\sup_{z\in\Gamma}
\big|\langle k_{h,\theta},(z-A_h(\theta))^{-1}p_h\rangle
-\langle k_{\theta},(z-A(\theta))^{-1}p\rangle\big|,
\]
and, using
\(\partial_\theta(z-A_h)^{-1}=-(z-A_h)^{-1}
\operatorname{diag}(k_{h,\mathrm{pat}}-\bar k_h)(z-A_h)^{-1}\)
(the path is affine,
\(A_h(\theta)=A_h(0)-\theta\operatorname{diag}(k_{h,\mathrm{pat}}-\bar k_h)\)),
the mixed defect \(\eta_h'\) as the same supremum for the two-term derivative
form
\(\langle k_{h,\mathrm{pat}}-\bar k_h,(z-A_h)^{-1}p_h\rangle
-\langle k_{h,\theta},(z-A_h)^{-1}
\operatorname{diag}(k_{h,\mathrm{pat}}-\bar k_h)(z-A_h)^{-1}p_h\rangle\)
against its continuum counterpart. Then, for \(j\le3\), \(l\in\{0,1\}\),
\(j+l\le3\),
\[
\sup_N\big|\partial_\theta^l\partial_t^j(f_h-f)\big|
\;\le\;C_\Gamma^{(j)}(t_1,t_2)\,\big(\eta_h+l\,\eta_h'\big),
\qquad
C_\Gamma^{(j)}
=\frac1{2\pi}\int_\Gamma\max(1,|z|)^{j}\,
\max\!\big(\e^{t_1\operatorname{Re}z},\e^{t_2\operatorname{Re}z}\big)
\,|\dd z|,
\]
finite because on the rays
\(\operatorname{Re}z=\omega-r\sin\delta'\), so the integrand decays like
\(r^{j}\e^{-t_1r\sin\delta'}\); the constant grows like
\(t_1^{-(j+1)}\) as \(t_1\downarrow0\), which is why the strictly positive
window is load-bearing. Hypotheses (H1), (H1$'$) and the decay of
\(\eta_h,\eta_h'\) are \emph{not} proved here for the upwind or
exponentially fitted schemes; they are the precise analytic content that the
observed ladder rates instantiate numerically.

\paragraph{Quantitative fold transfer.}
Let \(F,F_h\) have continuous \(\partial_t^jF\) (\(j\le3\)) and
\(\partial_\theta\partial_t^jF\) (\(j\le2\)) on \(N\). Assume \(F\) has an
interior fold at \(q_*=(t_*,\theta_*)\):
\(F_t(q_*)=F_{tt}(q_*)=0\), with margins
\(|F_{ttt}|\ge\mu_3>0\) and \(|F_{t\theta}|\ge\mu_\theta>0\) on \(N\).
Let \(\Lambda\) bound the sup norms
\(\|F_{t\theta}\|_N,\|F_{ttt}\|_N,\|F_{tt\theta}\|_N\), and let \(\omega\)
(part of the data) be a common nondecreasing modulus of continuity for
these three functions on \(N\); write
\(\hat\omega(r)=\max\big(\omega(r),\sqrt2\,\Lambda r\big)\) for the
augmented modulus that also controls the \(F_{tt}\) entry. Define
\[
\varepsilon_h=\max_{j=1,2}\|\partial_t^j(F_h-F)\|_{N},
\quad
\varepsilon_h'=\max\Big(\|\partial_t^3(F_h-F)\|_{N},
\max_{j\in\{1,2\}}\|\partial_\theta\partial_t^j(F_h-F)\|_{N}\Big),
\quad
\bar\varepsilon_h=\max(\varepsilon_h,\varepsilon_h').
\]
Let \(H=(F_t,F_{tt})\) with Jacobian
\(J=\begin{pmatrix}F_{tt}&F_{t\theta}\\F_{ttt}&F_{tt\theta}\end{pmatrix}\),
so \(\det J(q_*)=-F_{t\theta}(q_*)F_{ttt}(q_*)\) and
\[
\sigma:=1/\|J(q_*)^{-1}\|\;\ge\;
\frac{\mu_3\mu_\theta}{\sqrt{\mu_\theta^2+2\Lambda^2}}\,>0 .
\]
Let \(r_0=\min\big(r_\omega,\operatorname{dist}(q_*,\partial N)\big)\),
where \(r_\omega\) is any radius with \(\hat\omega(r_\omega)\le\sigma/16\)
(in particular \(r_\omega\le\sigma/(16\sqrt2\Lambda)\); so \(r_0\) is
determined by \((\mu_3,\mu_\theta,\Lambda,\omega)\) and the
position of \(q_*\) in \(N\)), and set
\[
\delta_0=\min\Big(\frac{\sigma}{32},\ \frac{\mu_3}{2},\
\frac{\mu_\theta}{2},\ \frac{5\,\sigma r_0}{8\sqrt2}\Big).
\]
If \(\bar\varepsilon_h<\delta_0\), then \(F_h\) has exactly one
solution \(q_h\) of \(F_{h,t}=F_{h,tt}=0\) in \(B(q_*,r_0)\); it is a
nondegenerate fold of \(F_h\), with margins
\(|F_{h,ttt}|\ge\mu_3-\bar\varepsilon_h>0\) and
\(|F_{h,t\theta}|\ge\mu_\theta-\bar\varepsilon_h>0\) on all of \(N\); and
\[
\|q_h-q_*\|\;\le\;2\sqrt2\,\|J(q_*)^{-1}\|\,\varepsilon_h .
\]
\emph{Converse (ladder version).} Suppose each \(F_h\) has a fold
\(q_h\in N\) with \(\operatorname{dist}(q_h,\partial N)\ge d>0\) uniformly,
the margins \(|F_{h,ttt}|\ge\mu_3\), \(|F_{h,t\theta}|\ge\mu_\theta\) hold
on \(N\), the functions
\(F_{h,t\theta},F_{h,ttt},F_{h,tt\theta}\) are bounded by \(\Lambda\) with
a common nondecreasing modulus \(\omega\) uniformly in \(h\), and the jets
\(\partial_t^jF_h\) (\(j\le3\)), \(\partial_\theta\partial_t^jF_h\)
(\(j\in\{1,2\}\)) converge uniformly on \(N\) to those of \(F\). Then for all
sufficiently fine \(h\) the limit \(F\) has a nondegenerate fold \(\bar q\)
with
\[
\|\bar q-q_h\|\le
2\sqrt2\,\frac{\sqrt{\mu_\theta^2+2\Lambda^2}}{\mu_3\mu_\theta}\,
\varepsilon_h ,
\]
and \(\bar q\) is the \emph{only} fold of \(F\) in \(N\): the \(N\)-wide
margins force \(F_{ttt}\) and \(F_{t\theta}\) to keep constant signs on the
connected rectangle, so \(F_t(\cdot,\theta)\) has a one-signed extremum
with value zero at any fold and strict monotonicity of \(F_t\) in
\(\theta\) excludes a second fold at any other \(\theta\). The proof (Supplemental
Material) is a simplified-Newton contraction seeded at the fold; no
Lipschitz control of the defect \(F_h-F\) is required.

\paragraph{Remark (a posteriori interpretation).}
Three statements of decreasing strength organize how the ladder is read.
(R1) \emph{Rigorous-conditional:} if the jets of \(F_h\) converge uniformly
on \(N\) to some \(F\) and the discrete folds, margins, and interior
distances satisfy the converse hypotheses uniformly along the ladder, then
\(F\) has a nondegenerate fold and every located discrete fold lies within
\(2\sqrt2\,\sqrt{\mu_\theta^2+2\Lambda^2}\,(\mu_3\mu_\theta)^{-1}
\varepsilon_h\) of it.
(R2) \emph{Falsification bound:} if two discretizations both satisfy the
transfer hypotheses toward the same limit and each located fold is the
unique one of the theorem's ball (both schemes track the same fold branch),
their located critical controls obey
\(|\theta_c^{(1)}(h)-\theta_c^{(2)}(h)|
\le2\sqrt2\,\|J(q_*)^{-1}\|(\varepsilon_h^{(1)}+\varepsilon_h^{(2)})\);
a large cross-scheme gap therefore falsifies joint convergence at the
claimed rates, while a small gap is necessary but not sufficient---error
components shared by both schemes (the cell-averaged masks and the budget
rule are common to both) cancel to leading order in the gap and are not
tested by it.
(R3) \emph{Error estimate under an asymptotic-regime assumption:} if the
level-to-level jet increments
\(d_k=\max_{j\le2}\|\partial_t^j(F_{h_{k+1}}-F_{h_k})\|_N\)
obey the geometric domination \(d_{m+1}\le\varrho\,d_m\) for \emph{all}
\(m\ge k\) with a fixed \(\varrho<1\)---the observed level ratios certify
this only on the computed levels, so \(\varrho\) should dominate, not
average, them, and the extension to uncomputed levels is an assumption---
then the telescoping bound
\(\tilde\varepsilon_{h_k}:=\max_{j\le2}\|\partial_t^j(F_{h_k}-F^{\rm lad}_\infty)\|_N
\le d_k/(1-\varrho)\)
controls the distance to the ladder's own jet limit \(F^{\rm lad}_\infty\);
its identification
with the continuum \(F\) is exactly the (H1)/(H1$'$) content. Richardson
extrapolants of \(\theta_c(h)\) are reported as estimates only: no
convergence rate is claimed for them unless the fitted leading coefficient
is nonzero, the observed order is stable across consecutive level triples,
and a strengthened remainder \(O(h^{p+s})\), \(s>0\), is assumed. The
cross-scheme gap of the two extrapolants is reported alongside them as the
honesty metric.

% ============ SM PROOF BLOCK ============

\section{Proof of the fold-transfer theorem}
\label{sm:fold-transfer-proof}

Write \(q=(t,\theta)\), \(H(q)=(F_t,F_{tt})(q)\),
\(H_h(q)=(F_{h,t},F_{h,tt})(q)\), \(J=DH\). The four entries of \(J\)
admit the augmented modulus \(\hat\omega(r)=\max(\omega(r),\sqrt2\Lambda r)\):
the entry \(F_{tt}\) is
\(\sqrt2\Lambda\)-Lipschitz (its gradient is \((F_{ttt},F_{tt\theta})\),
bounded entrywise by \(\Lambda\)), and \(F_{t\theta},F_{ttt},F_{tt\theta}\)
carry \(\omega\) by hypothesis, so each entry of \(J(q)-J(q')\) is at most
\(\hat\omega(\|q-q'\|)\), and \(\hat\omega(r_0)\le\hat\omega(r_\omega)\le
\sigma/16\) by the (nondecreasing) definition of \(r_\omega\). For a
\(2\times2\) matrix with entries bounded by \(\hat\omega\), the operator
norm is at most \(2\hat\omega\); we use this factor explicitly. All norms are
Euclidean and induced operator norms, and \(B(q_*,r_0)\subseteq N\) by the
definition of \(r_0\).

\emph{Step 0 (seed conditioning).} At \(q_*\), \(F_{tt}=0\),
\(|F_{ttt}|\ge\mu_3\), and \(|F_{ttt}|,|F_{tt\theta}|\le\Lambda\), so
\[
\sigma=\sigma_{\min}(J(q_*))
=\frac{|\det J(q_*)|}{\sigma_{\max}(J(q_*))}
\ge\frac{|F_{t\theta}|\,\mu_3}{\sqrt{F_{t\theta}^2+2\Lambda^2}}
\ge\frac{\mu_3\mu_\theta}{\sqrt{\mu_\theta^2+2\Lambda^2}},
\]
using \(\sigma_{\max}\le\|J\|_F=\sqrt{F_{t\theta}^2+F_{ttt}^2+F_{tt\theta}^2}\)
and monotonicity of \(x\mapsto x/\sqrt{x^2+2\Lambda^2}\).

\emph{Step 1 (perturbation inventory).} The defect matrix \(DH_h-DH\) has
entries bounded by \(\varepsilon_h\) (entry \(F_{tt}\)) and
\(\varepsilon_h'\) (entries \(F_{t\theta},F_{ttt},F_{tt\theta}\)), so
\(\|DH_h-DH\|\le\sqrt{\varepsilon_h^2+3\varepsilon_h'^2}
\le2\bar\varepsilon_h\). By \(\hat\omega(r_0)\le\sigma/16\),
\(\|J(q)-J(q_*)\|\le2\,\hat\omega(r_0)\le\sigma/8\) on \(B(q_*,r_0)\).

\emph{Step 2 (simplified Newton).} Assume
\(\bar\varepsilon_h<\delta_0\le\sigma/32\), so
\(2\bar\varepsilon_h\le\sigma/16\) and \(4\bar\varepsilon_h\le\sigma/8\).
At the seed,
\(\|DH_h(q_*)-J(q_*)\|\le2\bar\varepsilon_h\le\sigma/16\), so by the
Neumann series
\(\|DH_h(q_*)^{-1}\|\le\tfrac{16}{15}\|J(q_*)^{-1}\|
\le\tfrac87\|J(q_*)^{-1}\|\). Define
\(\Phi(q)=q-DH_h(q_*)^{-1}H_h(q)\). For \(q\in B(q_*,r_0)\),
\[
\|D\Phi(q)\|
=\|DH_h(q_*)^{-1}\big(DH_h(q_*)-DH_h(q)\big)\|
\le\frac87\cdot\frac1\sigma\Big(2\,\hat\omega(r_0)+4\bar\varepsilon_h\Big)
\le\frac87\cdot\frac1\sigma\cdot\frac{\sigma}{4}=\frac27 ,
\]
since
\(\|DH_h(q_*)-DH_h(q)\|\le\|(DH_h-DH)(q_*)\|+\|DH(q_*)-DH(q)\|
+\|(DH-DH_h)(q)\|\le2\bar\varepsilon_h+2\hat\omega(r_0)+2\bar\varepsilon_h\).
Also \(\|\Phi(q_*)-q_*\|=\|DH_h(q_*)^{-1}H_h(q_*)\|\le\alpha\) with
\(\alpha:=\tfrac{8\sqrt2}{7}\|J(q_*)^{-1}\|\varepsilon_h\), because
\(\|H_h(q_*)\|=\|(H_h-H)(q_*)\|\le\sqrt2\,\varepsilon_h\). The fourth term
of \(\delta_0\) gives
\(\tfrac75\alpha=\tfrac{8\sqrt2}{5\sigma}\varepsilon_h\le r_0\), so
\(\Phi\) maps \(B(q_*,\tfrac75\alpha)\subseteq B(q_*,r_0)\) into itself
(\(\alpha+\tfrac27\cdot\tfrac75\alpha=\tfrac75\alpha\)) and is a
\(\tfrac27\)-contraction there. The Banach fixed point \(q_h\) satisfies
\(H_h(q_h)=0\) and
\[
\|q_h-q_*\|\le\frac{\alpha}{1-\tfrac27}
=\frac75\alpha
=\frac{8\sqrt2}{5}\,\|J(q_*)^{-1}\|\,\varepsilon_h
\;\le\;2\sqrt2\,\|J(q_*)^{-1}\|\,\varepsilon_h .
\]

\emph{Step 3 (margins and uniqueness in the fixed ball).} The margin
hypotheses hold on all of \(N\), so directly
\(|F_{h,ttt}|\ge\mu_3-\bar\varepsilon_h>\mu_3/2>0\) and
\(|F_{h,t\theta}|\ge\mu_\theta-\bar\varepsilon_h>\mu_\theta/2>0\)
everywhere on \(N\), in particular at \(q_h\): the fold is nondegenerate
with the stated margins. If \(H_h(a)=H_h(b)=0\) with
\(a,b\in B(q_*,r_0)\), the averaged Jacobian
\(M=\int_0^1DH_h(a+s(b-a))\,\dd s\) obeys \(M(b-a)=0\) and
\[
\|M-J(q_*)\|\le2\,\hat\omega(r_0)+2\bar\varepsilon_h
\le\frac\sigma8+\frac\sigma{16}=\frac{3\sigma}{16}<\sigma,
\]
so \(M\) is invertible and \(a=b\).

\emph{Converse.} Under the ladder hypotheses, apply Steps 0--3 with the
roles of \(F\) and \(F_h\) exchanged, seeded at \(q_h\): Step 0 gives
\(\sigma_h=1/\|J_h(q_h)^{-1}\|\ge
\mu_3\mu_\theta/\sqrt{\mu_\theta^2+2\Lambda^2}\) uniformly in \(h\); the
common augmented modulus \(\hat\omega\) gives a common \(r_0\), and
\(\operatorname{dist}(q_h,\partial N)\ge d\) keeps
\(B(q_h,\min(r_0,d))\subseteq N\); uniform jet convergence makes the
defects \(\varepsilon_h,\bar\varepsilon_h\to0\), so the threshold
\(\delta_0\) (with \(r_0\) replaced by \(\min(r_0,d)\)) is met for fine
\(h\), and \(F\) acquires a unique nondegenerate fold \(\bar q\) in the
ball with
\(\|\bar q-q_h\|\le2\sqrt2\|J_h(q_h)^{-1}\|\varepsilon_h
\le2\sqrt2\,\sqrt{\mu_\theta^2+2\Lambda^2}\,(\mu_3\mu_\theta)^{-1}
\varepsilon_h\).\hfill\(\square\)
